\clearpage\begingroup\SourceLeaf{49c}\HistoricalHead{44}{L}{PHYSIK. BRIEFWECHSEL.}
\begin{center}
[9.]\\[.7em]
\textsc{Gauss} an \textsc{Alexander von Humboldt}. Göttingen, 10. Mai 1853.
\par\smallskip\rule{.72\textwidth}{.35pt}\\
{\small [Briefe zwischen A. v. \textsc{Humboldt} und \textsc{Gauss}, Leipzig 1877, S. 66.]}\\
\rule{.72\textwidth}{.35pt}
\end{center}
\noindent .\ .\ .\ .\ .\ .\ .\ .\ .\ .\ .\ .\ \hspace{.6em}.\ .\ .\ .\ .\ .\ .\ .
\par\noindent\hspace*{\parindent}In der letzten Zeit habe ich mich mit der Ausführung eines Apparates
beschäftigt, um die \textsc{Foucauld}schen Versuche in anderer Gestalt auszuführen.
Ich habe es bei diesen, so wie \textsc{Foucauld} selbst, \textsc{Secchi} u. a. sie ausgeführt
haben, wie einen grossen Mangel betrachtet, dass dazu ein Local erforderlich
wird, wie es an wenig Orten zu Gebote steht. \textsc{Secchi} hat, wenn ich nicht
irre, eine Höhe von mehr als 100 Fuss, \textsc{Foucauld} eines von mehr als 200,
\textsc{Garthe}[\textsuperscript{*})]\begingroup\renewcommand{\thefootnote}{}\footnotetext{[\textsuperscript{*})] \textsc{Caspar Garthe}, \textit{Foucaults Versuch als directer Beweis der Axendrehung der Erde, angestellt im Dom zu Cöln} u.s.w., Cöln 1852.]}\endgroup\ 134 u.s.w. Höhe. Mein Apparat ist in jedem Local anwendbar,
\clearpage\SourceLeaf{50}\HistoricalHead{45}{R}{GAUSS AN ALEXANDER V. HUMBOLDT.}
\noindent und zeigt schon jetzt die Einwirkung der Erdrotation nach kurzer Zeit auf
das schlagendste, ich hoffe aber (da er jetzt noch unvollständig ist) die noch
fehlenden Stücke, vielleicht successive, dahin zu bringen, dass alles in höchster
Eleganz und Präcision erscheint.\quad .\ .\ .\ .\ .\ .\ .\ .\ .\ .\ .\ .\ .\ .\ .\ .\ .\ .\ .\ .\ .

\bigskip\begin{center}\textbf{BEMERKUNG.}\end{center}
{\small\noindent\hspace*{\parindent}Das in den vorstehenden Briefstellen erwähnte von \textsc{Gauss} angegebene Pendel zur Ausführung des
\textsc{Foucault}schen Versuches ist erhalten; eine genauere Beschreibung findet man in dem Aufsatze »\textsc{Gauss} als
Physiker« in der zweiten Abteilung dieses Bandes.
\par\raggedleft\textsc{Schaefer}.\par}
\begin{center}\rule{.18\textwidth}{.35pt}\end{center}
\clearpage\endgroup
